\chapter{Отладка машины}
\chaptermark{Отладка машины}
\label{sec:debugging}

\section{Как отлаживают машину без схемы}

Инженер, которому дали работающую плату без схемы, отлаживает её четырьмя приёмами. Он ставит \emph{щуп} и смотрит, что идёт по проводу. Он \emph{подаёт тестовый сигнал} и смотрит, что изменилось. Он \emph{отключает} кусок схемы и смотрит, что перестало работать. И он \emph{читает регистры управления}, чтобы узнать, в каком режиме работает плата. Вся эта книга --- попытка прочитать схему мозга, и в этой главе посмотрим, какими из четырёх приёмов инженеры уже умеют пользоваться и что им говорят предыдущие главы о том, чего ждать.

Самый заметный сегодня пример такой отладки --- интерфейсы мозг--компьютер: устройства, которые читают импульсы нейронов и превращают их в команды для внешних машин, или, наоборот, подают в мозг сигналы извне. Им и посвящена большая часть главы, и в первую очередь тому, что делает компания Neuralink.

\section{Щуп: что слышит электрод}

Щуп мозга --- тонкий электрод, вставленный в ткань. Он слышит импульсы нейронов, лежащих в радиусе порядка сотни микрометров, --- обычно от одного до нескольких. Импульс на электроде --- всплеск в десятки--сотни микровольт длительностью около миллисекунды, и по его форме можно отличить соседние нейроны друг от друга. Лучше всего слышны крупные \nt{ГЛУТ}-нейроны: их в коре большинство (таблица~\ref{tab:types}), и их токи сильнее.

Сразу видна главная трудность. Хороший щуп сегодня --- это сотни или тысячи электродов, а нейронов в мозге $8.6\cdot10^{10}$. Мы подслушиваем примерно одну стомиллионную машины. Почему же из такой ничтожной выборки вообще что-то можно понять? Ответ дала глава~\ref{sec:code}. Соседние нейроны шумят согласованно, и усреднение по группе упирается в предел (утверждение~\ref{prop:pooling}): сотня нейронов несёт почти столько же, сколько тысяча. К тому же задача, которую решает двигательная кора, маломерна: у руки немного суставов (глава~\ref{sec:muscles}), и активность сотен нейронов, управляющих ею, укладывается в пространство всего из десятка-другого общих переменных~\cite{Gallego2017}. Щуп на сотню нейронов слышит почти все эти переменные. Если бы мозг хранил независимое сообщение в каждом нейроне, такое подслушивание было бы бесполезно.

\section{Поток данных: сжатие прямо на щупе}

Щуп выдаёт огромный поток. Чтобы увидеть форму импульса, каждый электрод нужно оцифровывать примерно $20$ тысяч раз в секунду с точностью около десяти бит. Тысяча электродов дают
\begin{empheq}[box=\keybox]{equation}
\begin{aligned}
R_{\text{сырой}} \;&=\; N\,f_s\,b \;\approx\; 10^3\cdot2\cdot10^4\cdot10 \;\approx\; 2\cdot10^8\ \text{бит/с},\\
R_{\text{имп}} \;&\approx\; N\,r\,\log_2\frac{e}{r\,\Delta t} \;\approx\; 8\cdot10^4\ \text{бит/с}.
\end{aligned}
\label{eq:probedata}
\end{empheq}
Вторая оценка --- ёмкость потока импульсов~\eqref{eq:spikeentropy} при $r=10$ импульсов в секунду и точности $\Delta t=1\ms$: всё, что в нём на самом деле есть, занимает в две с половиной тысячи раз меньше. Для вживлённого устройства это решающая разница: сырой поток нельзя передать по радио через кожу с тем питанием, которое можно позволить внутри головы, не нагревая ткань. Поэтому вживлённому устройству нужно выделять импульсы прямо на кристалле рядом с электродами и передавать наружу только события --- номер канала и момент импульса. В первом описании системы Neuralink до этого ещё не дошли: кристаллы усиливают и оцифровывают сигналы, весь сырой поток уходит по кабелю, а импульсы выделяет программа снаружи; сжатие на кристалле и связь без проводов названы там планом~\cite{Musk2019}.

Это знакомый приём. Глаз сжимает сигнал в сто раз, прежде чем отправить его по длинному кабелю (глава~\ref{sec:sensors}); событийная камера передаёт не кадры, а изменения. Щуп, чтобы уложиться в провод, вынужден делать то же, что делает сам мозг: говорить импульсами и только о новостях.

\section{Что делает Neuralink}

Устройство Neuralink --- щуп, собранный под эти ограничения~\cite{Musk2019}. Вместо жёсткой решётки иголок --- много гибких нитей тоньше волоса с электродами вдоль каждой: в первом описании системы --- до $3072$ электродов на $96$ нитях, по $32$ на нить, а в устройстве, вживлённом людям, по сообщениям компании, около тысячи каналов на $64$ нитях. Гибкие нити меньше повреждают ткань и лучше переносят то, что мозг внутри черепа всё время слегка сдвигается. Нити вставляет робот, обходя сосуды. В первом описании, как сказано выше, сырой поток~\eqref{eq:probedata} шёл наружу по кабелю. Устройство для людей, по сообщениям компании, устроено иначе: корпус полностью закрыт кожей, импульсы выделяются внутри, данные уходят по радио, питание приходит без проводов.

Задача первого устройства --- чтение. Нити ставят в ту часть коры, которая планирует движения руки, и декодер превращает импульсы в движение курсора на экране. Человек с параличом рук управляет компьютером, представляя движение. Сама идея не нова: интерфейсы на жёстких решётках из сотни электродов давно позволяли управлять курсором~\cite{Hochberg2006}, писать текст, представляя движения руки при письме~\cite{Willett2021}, и даже переводить попытки говорить в текст на экране со скоростью около шестидесяти слов в минуту~\cite{Willett2023}. Новое у Neuralink --- инженерия: число каналов, беспроводность, закрытый корпус и роботизированная вставка, то есть попытка сделать щуп, который можно носить годами вне лаборатории. Насколько нам известно, результаты испытаний на людях публиковались в основном самой компанией, а не в рецензируемых статьях; компания сообщала, в частности, что часть нитей после вставки сместилась и число работающих каналов уменьшилось. Для отладки это обычная неприятность: щуп, который двигается относительно платы, слышит уже другие провода.

Второе направление компании --- запись: устройство для зрения, которое подаёт сигналы в зрительную кору человека, потерявшего зрение. О нём ниже, в разделе о записи.

\section{Чтение: декодер как получатель}

Декодер интерфейса --- это получатель в смысле утверждения~\ref{prop:reader}: он сам решает, какую часть сообщения прочитать. Практически все интерфейсы читают \emph{частоты групп}: считают импульсы каждого канала в окнах по несколько десятков миллисекунд и складывают их с весами, подобранными так, чтобы получалось задуманное движение. Это частотный код группы из раздела~\ref{sec:rate}, и он выбран не случайно: при нескольких импульсах на окно частота одного нейрона не определена, а сумма по сотне уже работает.

Сколько информации удаётся вынуть? Письмо «мысленной рукой» шло со скоростью около девяноста знаков в минуту~\cite{Willett2021}, то есть полтора знака в секунду: даже при пяти битах на знак это меньше восьми бит в секунду. Управление курсором, по сообщениям Neuralink, даёт порядка десяти бит в секунду. А верхняя граница~\eqref{eq:spikeentropy} для \emph{одного} нейрона --- десятки бит в секунду, для сотни --- тысячи. Интерфейс вынимает из подслушанных нейронов малую долю того, что они могли бы нести. Узкое место --- не провод, а то, сколько независимых переменных несёт группа (утверждение~\ref{prop:pooling}) и насколько хорошо декодер понимает код, который, как мы видели в главе~\ref{sec:code}, до конца не расшифрован.

Есть и вторая особенность, знакомая по главе~\ref{sec:motorlearning}. Декодер и мозг учатся друг у друга. Мозг видит, куда поехал курсор, получает ошибку и подстраивает свои команды --- то же обучение движениям по проводу ошибки. А инженеры время от времени заново подбирают веса декодера, потому что нейроны под нитями меняются и связи в сети переучиваются. Получаются два ученика, которые подстраиваются друг под друга; чтобы такая пара не ушла вразнос, скорости их обучения удобно развести: пусть один учится быстро, а другой медленно.

\begin{keyblock}
\begin{proposition}[Почему щуп на тысячу нейронов работает]
\label{prop:probe}
Интерфейс, подслушивающий ничтожную долю нейронов, может управлять движением потому, что активность группы маломерна и её шумы согласованы: несколько сотен нейронов несут почти все общие переменные. По той же причине интерфейс вынимает лишь несколько бит в секунду --- столько, сколько независимых переменных в задаче, а не столько, сколько вмещает поток импульсов. Работа интерфейсов измерена; насколько лучше станет декодер, когда код будет понят, --- открытый вопрос.
\end{proposition}
\end{keyblock}

\section{Запись: труднее, чем чтение}

Подать сигнал в машину труднее, чем прочитать. Электрод, через который пропускают ток, возбуждает не один нейрон, а все нейроны и провода вокруг себя, без разбора типа и знака. Записать им код, в котором важно, \emph{какой} нейрон и \emph{когда} сработал (глава~\ref{sec:code}), нельзя: можно лишь грубо задать, \emph{где} и \emph{насколько сильно}.

Для зрения этого отчасти хватает. Ток через электрод в зрительной коре человек видит как светящуюся точку в определённом месте поля зрения; когда точки зажигали одновременно, до шестнадцати сразу, человек, потерявший зрение, узнавал некоторые буквы и различал границы предметов~\cite{Fernandez2021}. Это код местом, и его записать можно. Но вспомните главу~\ref{sec:sensors}: глаз посылает в мозг не пиксели, а сжатое описание --- края, перемены, посветление и потемнение, по отдельным проводам. Запись «точек света» в кору --- это запись пикселей в машину, которая ждёт на входе совсем другой код. Поэтому искусственное зрение пока грубее, чем можно было бы ждать по числу электродов. Есть и испытательные сигналы, которым электрод не нужен: зрительные иллюзии подают через глаз необычный вход и выводят машину из обычного режима, а каждая ошибка указывает на свой механизм (раздел~\ref{sec:illusions}).

\section{Чего щуп не видит}

Электрод слышит адресную шину --- импульсы \nt{ГЛУТ}-нейронов и хуже, мелких \nt{ГАМК}. Но в главах~\ref{sec:serocomputer}--\ref{sec:acet} мы видели, что режим работы всей машины задают широковещательные регистры: концентрации \nt{СЕРО}, \nt{ДОФА}, \nt{НОРА}, \nt{АЦЕТ}. Их электрод не слышит: нейронов-источников ничтожно мало, и их сигнал --- не импульс рядом с щупом, а концентрация в объёме. Отладчик, который видит шину данных, но не видит регистров управления, всегда будет удивляться: один и тот же вход будет давать разные ответы, потому что где-то изменился $g$, $a$ или $\tau_\gamma$. Концентрации можно мерить химическими датчиками прямо в ткани~\cite{Bunin1998}, но в интерфейсах такие датчики пока не используются. Совсем вне поля зрения щупа остаётся и второй, медленный выход машины --- гормоны в крови (глава~\ref{sec:chemoutput}): их меряют в пробах крови, а не электродом.

Не видит щуп и весов --- а именно в них хранится почти вся память и всё выученное. Их можно лишь угадывать по тому, как меняются ответы. И не видит он боли так, как её чувствует человек: в главе~\ref{sec:pain} мы видели, что сила сигнала тревоги задаётся самой машиной --- воротами в спинном мозге и регуляторами сверху, так что по датчику нельзя прочитать, насколько больно.

\section{Итог: отладка и открытые вопросы}

\begin{table}[t]
\centering
\footnotesize
\setlength{\tabcolsep}{4pt}
\renewcommand{\arraystretch}{1.15}
\begin{tabular}{>{\raggedright}p{2.2cm}>{\raggedright}p{3.6cm}>{\raggedright}p{3.4cm}>{\raggedright\arraybackslash}p{4.0cm}}
\toprule
Приём отладки & Что делает интерфейс & Что объясняет книга & Что известно \\
\midrule
щуп & слышит сотни--тысячи нейронов & маломерность и согласованный шум (гл.~\ref{sec:code}) & измерено \\
сжатие на щупе & передаёт события вместо сырого сигнала & ёмкость потока импульсов~\eqref{eq:spikeentropy} & инженерно решено \\
чтение & частоты групп $\to$ движение, текст, речь & декодер --- получатель, частотный код группы & работает; несколько бит в секунду \\
совместное обучение & мозг и декодер подстраиваются друг под друга & провод ошибки (гл.~\ref{sec:motorlearning}) & наблюдается; правила подстройки --- предмет исследований \\
запись & ток зажигает точки в поле зрения & код местом записать можно, временной --- нет (гл.~\ref{sec:sensors}) & точки измерены; полезное зрение --- пока нет \\
чтение регистров & почти не делается & широковещательные каналы (гл.~\ref{sec:serocomputer}--\ref{sec:acet}) & химические датчики есть в опытах \\
\bottomrule
\end{tabular}
\caption{Отладка машины: что умеют интерфейсы мозг--компьютер и что о них говорят главы книги.}
\label{tab:debugging}
\end{table}

Таблица~\ref{tab:debugging} сводит главу. Интерфейсы мозг--компьютер уже умеют ставить щуп на адресную шину и читать с неё несколько бит в секунду, и то, что это вообще работает, объясняется маломерностью и согласованным шумом из главы~\ref{sec:code}. Записывать в машину они умеют лишь грубо, потому что её входной код сжат и адресован отдельным проводам. А регистры управления и веса --- то, что делает машину обучаемой и настраиваемой, --- щупу пока почти не видны.

Каждый открытый вопрос главы~\ref{sec:synthesis} --- это и задача для отладчика. Какой код читать, решится, когда щупы станут плотнее и декодеры научатся пользоваться временем импульсов, а не только частотами. Как учится многослойная сеть, можно проверить, поставив щупы сразу в несколько слоёв и глядя, как меняются их ответы при обучении. Что передают широковещательные каналы, станет видно, когда к электрическим щупам добавятся химические. Эта книга --- попытка прочитать схему машины по её поведению и по законам, которые знает любой инженер. Отладчики, которые строятся сейчас, дадут возможность проверить эту схему щупом.

\section{Задачи}

\begin{exercise}[\lvlA]
\label{ex:17:1}
\emph{Бюджет радиоканала.} (а)~Пересчитайте~\eqref{eq:probedata} и отношение двух потоков для $N=1000$ и $N=3072$ электродов.
(б)~Пусть передача через кожу стоит $1$~нДж на бит, а на радиопередатчик можно тратить не больше $10$~мВт. Какая мощность нужна для сырого потока и для потока импульсов? Какой энергии на бит потребовал бы сырой поток?
(в)~Импульсы передаются кадрами по $1\ms$: заголовок $16$~бит и $10$-битный номер канала на каждый импульс. Найдите средний поток при $N=1000$, $r=10$~имп/с и сравните его с ёмкостью~\eqref{eq:probedata}.
\end{exercise}

\begin{exercise}[\lvlA]
\label{ex:17:2}
\emph{Сколько бит в сотне нейронов.} Декодер складывает импульсы $N$ нейронов за окно $T=50\ms$, каждый работает с частотой около $r=20$~имп/с, шум пуассоновский.
(а)~По~\eqref{eq:rateerror} и~\eqref{eq:correlated} найдите относительную погрешность суммы при $N=1$, $10$, $100$, $1000$ для $c=0$ и $c=0.1$.
(б)~Пусть управляемая переменная меняет частоту группы со среднеквадратичной амплитудой $30\%$, а шум нормальный. Оцените информацию на окно $\tfrac12\log_2(1+\text{SNR})$ и верхнюю границу скорости, считая окна независимыми.
(в)~Сравните с числами из текста. Какое предположение о коррелированном шуме делает эту оценку пессимистичной?
\end{exercise}

\begin{exercise}[\lvlB]
\label{ex:17:3}
\emph{Скорость интерфейса-клавиатуры.} Интерфейс выбирает один из $N$ знаков; с вероятностью $P$ выбран задуманный, ошибки равновероятны среди остальных $N-1$.
(а)~Докажите, что взаимная информация на выбор при равновероятных знаках равна
$B=\log_2N+P\log_2P+(1-P)\log_2\dfrac{1-P}{N-1}$.
(б)~Докажите, что это ёмкость такого канала.
(в)~Найдите $B$ и скорость в битах в секунду при $N=32$, полутора знаках в секунду и $P=0.99$, $0.94$, $0.9$, $0.8$. Сверьте с оценкой из текста.
(г)~При каком $P$ скорость нулевая? Сколько знаков в секунду нужно при $P=0.94$, чтобы получить $10$~бит/с?
\end{exercise}

\begin{exercise}[\lvlB]
\label{ex:17:4}
\emph{Моделирование: декодер группы.} $N$ нейронов настроены на двумерную скорость $\mathbf v$: частота $\lambda_i=\bigl[b+m\,\mathbf v\cdot\mathbf p_i\bigr]_+$ с $b=20$~имп/с, $m=15$~имп/с, предпочтительные направления $\mathbf p_i$ равномерно по кругу. Каждая компонента $\mathbf v$ --- случайный процесс Орнштейна--Уленбека с постоянной $0.5$~с и стандартным отклонением $0.5$. Импульсы пуассоновские, окна $T=50\ms$, $2000$ окон. Чтобы смоделировать шум, похожий на сигнал, все нейроны получают вместо $\mathbf v$ одно и то же $\mathbf v+\boldsymbol\xi$, где $\boldsymbol\xi$ --- независимый в каждом окне нормальный шум со стандартным отклонением $\sigma_\xi$ по каждой компоненте.
(а)~Декодируйте вектором группы $\hat{\mathbf v}=\frac{2}{NmT}\sum_i(n_i-bT)\,\mathbf p_i$. Выведите дисперсию ошибки: $2b/(Nm^2T)+\sigma_\xi^2$ на компоненту.
(б)~Постройте среднеквадратичную ошибку от $N=10$, $30$, $100$, $300$, $1000$ для $\sigma_\xi=0$ и $0.2$ и сравните с формулой.
(в)~При каком $N$ два вклада в ошибку равны? Сколько бит на компоненту за окно даёт бесконечная группа при $\sigma_\xi=0.2$? Свяжите результат с утверждениями~\ref{prop:pooling} и~\ref{prop:probe}.
\end{exercise}

\begin{exercise}[\lvlB]
\label{ex:17:5}
\emph{Два ученика.} Мозг выдаёт команду $m=b\,v^*$ на задуманную скорость $v^*$ ($\langle v^{*2}\rangle=1$), декодер выдаёт $v=d\,m$. Ошибка $e=v-v^*$. Оба учатся градиентным спуском по $\tfrac12e^2$ со скоростями $\eta_b$ и $\eta_d$.
(а)~Запишите уравнения в непрерывном времени. Покажите, что величина $\eta_d b^2-\eta_b d^2$ сохраняется, и найдите, куда придёт пара, выйдя из $b=1.5$, $d=1$, при $\eta_b=1$, $\eta_d=0.1$ и наоборот.
(б)~Для одновременных обновлений после каждой попытки (оба ученика используют прежние значения $b$ и $d$) линеаризуйте ошибку $\varepsilon=db-1$ вблизи $db=1$ и докажите, что обучение сходится при $\eta_bd^2+\eta_db^2<2$.
(в)~Смоделируйте при $b_0=1.2$, $d_0=1$ и $\eta_b=\eta_d=0.8$, $0.9$, $1.1$, $1.2$, $1.5$, $2.0$. Что происходит, когда каждый ученик по отдельности устойчив, а вместе --- нет?
(г)~Сформулируйте правило выбора скоростей для инженера, который подстраивает декодер.
\end{exercise}

\begin{exercise}[\lvlB]
\label{ex:17:6}
\emph{Проектирование: конвейер данных вживлённого щупа.} Дано: $N=1024$ канала, $20$~тыс. отсчётов в секунду по $10$~бит, средняя частота нейронов $10$~имп/с, радио $1$~нДж/бит с бюджетом $1$~мВт, задержка до декодера не больше $25\ms$, вероятность потери кадра не больше $10^{-6}$. Шум отсчётов белый нормальный с дисперсией $\sigma^2$.
(а)~Выберите порог выделения импульса так, чтобы ложные срабатывания составляли не больше $1\%$ настоящих ($0.1$~Гц на канал). Что будет при пороге $3.5\sigma$?
(б)~Вариант А: кадры по $1\ms$ фиксированной длины на $C$ событий по $10$ бит плюс заголовок $16$ бит. Найдите наименьшее $C$ и мощность передатчика.
(в)~Вариант Б: числа импульсов каждого канала за $20\ms$, по $2$ бита с насыщением на $3$. Найдите поток, мощность и долю насыщенных отсчётов; оцените выигрыш от энтропийного кодирования.
(г)~Сравните варианты с сырым потоком, выберите решение и объясните, что сломается в каждом варианте при синхронной вспышке всей группы.
\end{exercise}

\begin{exercise}[\lvlC]
\label{ex:17:7}
\emph{Предел скорости интерфейса.} Утверждение~\ref{prop:probe} говорит, что скорость ограничена числом независимых переменных в задаче, а не ёмкостью потока импульсов.
(а)~Постройте верхнюю оценку скорости через размерность $D$ общих переменных, полосу $W$ их изменений и отношение сигнал/шум на переменную, $R\le DW\log_2(1+\text{SNR})$. Обоснуйте выбор чисел ссылками на главы книги и оцените $R$ при шуме, ограниченном корреляцией, и при независимом шуме.
(б)~Сравните с измеренными несколькими битами в секунду. Какие факторы объясняют разрыв?
(в)~Открытая часть. Предложите опыт на реальном интерфейсе, который различит два объяснения узкого места: шум, похожий на сигнал (утверждение~\ref{prop:pooling}), и незнание кода, в частности времён импульсов (глава~\ref{sec:code}). Что вы ожидаете увидеть на графике «извлечённая информация --- число используемых каналов» в каждом случае?
\end{exercise}
